\documentclass[12pt, a4paper]{article}
\usepackage[spanish]{babel}
\usepackage{color,graphicx}
\usepackage{graphics}
\usepackage{amsmath,amssymb,amsfonts,amsthm}
\usepackage[T1]{fontenc}
\usepackage{float}
\usepackage{xcolor}
\usepackage{amsthm}
\usepackage{tikz}
\usepackage{enumerate}
\usepackage{picinpar}
\usepackage{newtxtext,fourier}
\makeatother
\usepackage{thmtools}
\usepackage[framemethod=TikZ]{mdframed}
\mdfsetup{skipabove=1em,skipbelow=1em}

\newcommand{\sqed}{\hfill\includegraphics[height=2ex]{sancoya.png}}

\theoremstyle{definition}
\declaretheoremstyle[
    headfont=\bfseries\color{black!70!black}, bodyfont=\normalfont,
    mdframed={
        linewidth=0.7pt,
        rightline=true, topline=true, bottomline=true,
        linecolor=black, backgroundcolor=black!2!white,
    }
]{thmbox}

\declaretheorem[style=thmbox, name=Definición]{definition}
\declaretheorem[style=thmbox, numbered=no, name=Ejemplo]{example}
\declaretheorem[style=thmbox, name=Proposición]{prop}
\declaretheorem[style=thmbox, name=Teorema, numbered=no]{theorem}
\declaretheorem[style=thmbox, name=Lema]{lemma}
\declaretheorem[style=thmbox, numbered=no, name=Corolario]{corollary}
\declaretheorem[style=thmbox, numbered=no, name=Nota]{note}




\usepackage{geometry}
 \geometry{
 a4paper,
 total={170mm,260mm},
 left=25mm,
 top=30mm,
 }
\begin{document}
\setlength{\parindent}{0cm}
\hoffset-0.46cm
\voffset-1.46cm
\noindent
\hbox to \linewidth{
  \begin{minipage}[t]{40pt}
    \vspace*{0pt}
\includegraphics[scale=0.22]{Graphics/logo.eps}
  \end{minipage}
  \hss
  \begin{minipage}[t]{0.8\linewidth}
    \vspace*{15pt}\centering
     Universidad Nacional de Colombia\\
    \small Facultad de Ciencias\\
    Teoría analítica de números\\
    Tarea III
  \end{minipage}
  \hss
}
\phantom{xd}\\
\small
\textsf{Mateo Andrés Manosalva Amaris}
\normalsize
\dotfill
\vspace{0.3cm}

\textbf{Ejercicio 1.} Demuestre que la serie
$$\sum_p\dfrac{1}{p\log p}$$ es convergente

\begin{proof}
  Note que,

  $$\sum_{p\leq x}\frac{1}{p\log p}=\sum_{p\leq x}\frac{\log p}{p\log^2p},$$
  podemos convertir esta suma en una integral de Riemann-Stieltjess como sigue
  \begin{align*}
    \sum_{p\leq x}\frac{1}{p\log p}&=\int_{2^-}^x\frac{1}{\log^2t}d\left(\sum_{n\leq t}\frac{\mathbb{1}_p(n)\log(n)}{n}\right)\\
    &=\int_{2^-}^x\frac{1}{\log^2t}d\left(\log t+O(1)\right)\\
    &=\int_{2^-}^x\frac{1}{t\log^2t}dt\\
    &=\dfrac{1}{\log 2}-\frac{1}{\log x}\\
    &=O(1)
  .\end{align*}
\end{proof}
Luego tomando $x\to \infty$, acabamos
\end{document}